\subsection*{Drill: Hip Drive}
\label{drill:drives/hip}

This drill requires pairs of skaters. 

Starting engaged, one skater is driving the other blocking, with both skaters starting in the middle of the track. 

The driving skater should attempt to use their thigh and hip to push the other skater off the track. 
The skater being driven should attempt to hold their position (rather than simply skating away).


The core idea for the pushing skater is to collapse the leg of the other blocker - once their knee is pointing towards the edge of the track it is much easier to push them in that direction. 
To achieve this the block must begin legally - with either the hip or the thigh making contact. The skater can then apply pressure through the thigh to turn the blocker's leg.
If the drive is initiated with the thigh, then the hip can be swung in as an additional hit during the drive.

Practice in pairs, and rotate skaters where possible - driving is best practiced against a wide range of different opponents.
Also don't forget to practice driving to both the inside and outside.

{\it Progressions: }
A natural progression to just hip and thigh engagement is to blur the line with a shoulder drive - engaging that too}

% TODO: Figure
